\section{Introduction}
    % 赋予人形机器人类似人类的通用运动能力，是具身智能的终极愿景之一。得益于现代通用运动追踪器的飞速发展，机器人能够通过统一的底层控制器复现人类全身参考动作，而不必重新训练

	% 随着通用追踪器的能力提升，限制机器人运动能力提升的瓶颈逐渐从“能否控制机器人完成动作”转向“能否提供足够多样的参考动作”

	% 传统动作捕捉 (MoCap) 数据具有良好物理一致性和清晰标定，但采集成本高。相比之下，互联网视频覆盖了更丰富的极限动作和交互动作，是更具扩展性的参考动作来源

	% 但即便有规模优势，使用互联网视频依然潜伏着新的问题：我们不再执行经过标定、接触关系相对可靠的 MoCap 轨迹，而是执行由视觉系统从非受控视频中恢复出的运动参考

    % ---------

    % 然而，直接从野外视频中恢复出的三维人体运动轨迹，并不是为闭环物理控制器准备的参考轨迹

	% 单目估计、人体重建和机器人重定向会引入根节点高度漂移、足端穿地或漂浮、脚步滑移、朝向抖动以及接触时序不一致等伪影

	% 我们将这种“运动学上合理”到“动力学上可执行”之间的mismatch称为Kinematic-to-Dynamic Gap: 由视频、MoCap 或重定向流程得到的参考动作在运动学层面可以保持合理的人体姿态、关节轨迹和外观连续性，但它并不保证满足机器人闭环执行所需的动力学条件

	% 若迫使优化的底层模型强行追踪这些违背物理定律的运动学信号，必将诱发剧烈的动力学冲突，导致部署的灾难性失效

    % ---------

    % 这种失配并不只是定性的直觉；它可以直接反映在 tracker 面对扰动时还剩多少稳定裕度上

	% 如 Fig.~\ref{fig:robustness-budget} 所示，随着参考帧根节点噪声增大，机器人在推力后的最小稳定裕度持续下降，整段动作的 end-to-end success rate 也随之降低

    % ---------

    % 现有运动跟踪器已经能够复现复杂的人形动作。然而，这些方法通常假设参考动作本身是动力学可执行
	
	% 这一假设在视频重建动作中并不总是成立。视频重建动作通常在姿态空间中看起来自然。关节姿态可能连续，人体形态可能合理，根节点轨迹也可能平滑
	
	% 但是，这些动作仍然可能包含细小的参考帧误差。根节点高度、身体朝向、接触时序和水平位移的偏差都会影响物理执行
	
	% 我们将这种现象称为 Kinematic-to-Dynamic Gap：一个参考动作可以在运动学指标上接近目标动作，却仍然在动力学系统中难以被稳定执行。这个 gap 的存在意味着，跟踪失败并不一定来自跟踪策略本身能力不足，也可能来自上游参考动作已经消耗了跟踪器原本用于抵抗外界扰动和模型误差的鲁棒性预算。

    % ---------

    % 一种直接的解决思路是继续增强底层运动跟踪器，使其能够容忍更大范围的参考误差。然而，这种做法会将两个本应区分的问题混合在一起：一方面是学习如何控制类人机器人稳定执行给定动作，另一方面是修复由上游动作生成过程引入的参考帧伪影。
	
	% 前者属于动力学控制问题，后者则更接近参考动作的可执行性修正问题。如果只依赖重新训练或增强跟踪器，那么每当上游动作来源、视觉重建模型或误差分布发生变化时，都可能需要重新调整控制策略。
	
	% 这不仅增加训练成本，也会削弱方法的模块化和可迁移性。更重要的是，跟踪器本身已经学习到了一套复杂的低层控制结构；让它同时承担参考修复和动力学控制两种职责，可能会掩盖真正的问题来源。
	
	% 因此，我们采取另一种视角：与其替换或重训跟踪器，不如在参考动作进入跟踪器之前，对其进行轻量级的前端修正，使其在保持原始动作语义的同时，更接近下游跟踪器能够稳定执行的区域。

    % ---------

    % 基于这一思想，我们提出 FEMR（Front-End Motion Refiner），一个放置在冻结运动跟踪器之前的轻量级前端动作修正模块。
	
	% FEMR 的目标不是学习一个新的跟踪器，也不是在关节空间中重新生成完整动作，而是针对参考帧中的关键伪影预测小幅、任务空间的残差修正。
	
	% 具体而言，FEMR 接收当前跟踪观测以及近期 anchor 误差历史，输出作用于根节点平移和姿态的 $\Delta SE(3)$ 修正，并通过置信度门控决定何时介入、介入哪些维度以及修正幅度应当多大。
	
	% 这种设计使 FEMR 能够默认保持克制，在参考动作已经可执行时选择不修正；而当参考帧误差落入可修复边界时，则输出足够有效的修正以恢复下游跟踪器的执行余量。
	
	% 为了避免 FEMR 学成另一个隐式控制器，我们冻结原始 GMT 跟踪器，并将训练信号集中在修正前后参考动作的可执行性差异上。
	
	% 这样，FEMR 被明确地约束为一个前端动作修正器：它只负责减少上游参考误差对动力学执行的损害，而不改变底层跟踪器已经学习到的控制结构。